\section{音楽嗜好の共有が適している環境}\label{音楽嗜好の共有が適している環境}
% 〇〇ではこういうのが問題で～
音楽嗜好の共有を通じた自己開示には，音楽ストリーミングサービスでのプレイリストの共有，SNSでの音楽体験の発信，あるいは日常会話での好みの表明など，様々な形態が存在する．しかし，これらの方法では好みの表明が一方向的になりやすく，相互の深い理解や受容を促進する機会として十分に機能しない．

% こういう条件が必要で～
榎本が指摘するように，自己開示は状況と密接な関係があり，環境要因が大きく影響する．深田は，特に相手からの受容が自己開示を促進する重要な要因であることを示している．これらの知見に基づけば，効果的な自己開示の実現には，自己開示の相互性と受容を考慮した適切な環境が必要となる．そのため，音楽嗜好の共有を通じた自己開示には以下に挙げる3つの条件を満たす環境が望ましい．


\subsubsection{音楽の実体験の共有}
音楽嗜好の共有を通じた自己開示の場においては，単なる好みの表明だけでなく，実際の音楽を共に聴く，歌うといった形で体験できることが重要である．Tarrの研究によれば，音楽の共同体験は社会的絆の形成を促進し，他者との心理的な結びつきを強める効果がある．また，榎本が指摘するように，自己開示は言語的なものに限らず，非言語的な表現も含む多様な形態を持つ．音楽体験を共有することは，感情や印象をリアルタイムで共有することを可能にし，自己開示と自己開示に対する受容を促すことができる．


\subsubsection{明示的な相互性の確保}
\ref{状況要因}項でも述べたように，自己開示の状況要因で最も影響の強いとされるのは自己開示の相互性である．そのため，片方の一方的な開示に偏ってしまわないよう，音楽を聴く側と開示（表現）する側の役割が明確に分かれ，それらを交互に行う構造が必要である．


\subsubsection{対面での交流}
Weberの研究によれば，思いやりや共感の表現は相手に理解されているという感覚を高め，信頼関係を深める．対面での交流は，相手の反応や表情から受容を直接的に確認することを可能にし，それが次の自己開示を促進する重要な要因となる．また，Reisが指摘するように，対面でのコミュニケーションは相互理解の深化において重要な役割を果たす．

\vskip\baselineskip

% それを満たすのはカラオケで～
以上の条件を総合的に満たす環境として，本研究ではカラオケに着目する．カラオケは，音楽を通じた自己開示や自己表現が自然に行われる場\cite{烏賀陽弘道_2005}であるように，選曲や歌唱を通じた音楽の実体験を共有できる場と言える．さらに「歌う」「聴く」という明確な「開示者」と「被開示者」としての役割が存在することにより相互作用が生まれやすい．これらの特徴は，前述した条件を効果的に満たしており，音楽の共有を通じた自己開示をする環境として適しているといえる．